Experiment~3 showed that the epoch-aware proposed pipeline outperformed both comparison methods across the two driving corpora. Proposed-approach reached macro-$F_1$ of 88.25\% on Internal and 80.68\% on Cao2018, compared with 76.54\% and 69.40\% for BLINKER-concat and 66.64\% and 56.27\% for MNE-annot on the same two corpora (Table~\ref{tab:f1_significance}). BLINKER-concat represented an established blink-specific method applied to the concatenated recording~\citep{kleifges2017blinker}, whereas MNE-annot represented the default amplitude-annotation routine from MNE-Python.

The corresponding session-level differences were statistically significant for both BLINKER-concat and MNE-annot on both corpora, with $p<0.05$ after Bonferroni correction in every comparison, and neither baseline equalled or exceeded Proposed-approach in any channel group on either dataset. These comparisons establish the advantage of the complete proposed configuration under the common evaluation protocol. Because the experiment compared complete detection strategies rather than a screening-on versus screening-off ablation of an otherwise identical detector, it does not by itself attribute the improvement exclusively to epoch screening.
